% !TEX program = xelatex
\documentclass{misisarticle}

\grnti{06.35.51}
\asjc{2002}
\oecd{5.2}

\title{Математическое моделирование выбора потребителя между покупкой и подпиской с использованием теории принятия решений}
\misisauthor{А.~Д.~Картаева}{НИТУ «МИСИС»}{Москва, Россия}{}

\abstract{%
Рассматривается выбор между разовой покупкой и ежемесячной подпиской при неопределённом сроке использования продукта.
Модель учитывает приведённую стоимость платежей, остаточную стоимость покупки и отношение потребителя к риску.
Для известного срока использования выведен порог окупаемости покупки.
При случайном сроке альтернативы сравниваются по ожидаемой полезности.
На условных данных показано, как вероятность длительного использования и неприятие риска меняют предпочтительный вариант.%
}
\keywords{потребительский выбор, подписка, покупка, ожидаемая полезность, дисконтирование, срок использования}

\begin{document}
\maketitle

\section{Введение}

Один и тот же цифровой продукт может продаваться по бессрочной лицензии или предоставляться за ежемесячную плату. При покупке потребитель платит всю сумму в начале, а при подписке распределяет платежи во времени и может прекратить их при отказе от продукта. В момент выбора потребитель обычно располагает лишь оценкой будущего срока использования. Поэтому сравнение тарифов требует учитывать размер платежей и вероятность того, что доступ понадобится надолго.

Цель работы состоит в определении условий, при которых покупка предпочтительнее подписки. Вклад работы заключается в выводе порогов по длительности использования и абонентской плате в рамках одной модели и в оценке чувствительности решения к вероятности длительного использования и отношению к риску.

\section{Анализ работ}

Bardhi и Eckhardt~\cite{bardhi} исследовали потребление на основе доступа на примере каршеринга. Для настоящей работы важно различие между доступом и владением. В модели оно отражено в возможности прекратить платежи по подписке и в возможности получить остаточную стоимость покупки, если её можно перепродать.

Lambrecht и Skiera~\cite{lambrecht} выявили отклонения от выбора тарифа с минимальной фактической стоимостью при сравнении фиксированной оплаты и оплаты по объёму использования. DellaVigna и Malmendier~\cite{dellavigna} показали на данных фитнес-клубов, что участники с ежемесячным абонементом в среднем платили за посещение больше, чем могли бы при оплате пакета посещений. Эти результаты показывают, почему при анализе фактического выбора следует учитывать ожидания потребителя на момент принятия решения. Методологической основой сравнения альтернатив при заданных вероятностях служит теория ожидаемой полезности~\cite{keeney}.

\section{Постановка задачи}

Потребитель выбирает способ доступа к одному продукту. Покупка требует платежа $P>0$ в момент выбора. Подписка предусматривает плату $s>0$ в конце каждого месяца доступа и прекращается после последнего месяца использования без штрафа. Обе альтернативы дают одинаковую функциональность, тарифы постоянны, а потребитель может оплатить покупку сразу.

Пусть $T$~--- число месяцев использования, $r\geq0$~--- месячная ставка дисконтирования, а $R(t)$~--- выручка от продажи покупки в конце месяца $t$. Для непередаваемой лицензии $R(t)=0$. Предполагается, что срок использования продукта одинаков для обеих схем оплаты. При известном сроке $T=t$ сравниваются приведённые затраты. При случайном сроке заданы сценарии $(t_i,p_i)$, где $t_i\geq1$, $p_i=\Pr(T=t_i)$ и $\sum_i p_i=1$. Тогда сравнивается ожидаемая полезность альтернатив.

\section{Описание методов}

\subsection{Приведение затрат к моменту выбора}

Приведённая стоимость серии из $t$ единичных ежемесячных платежей равна
\begin{equation}
  A(t,r)=\sum_{k=1}^{t}(1+r)^{-k}=
  \begin{cases}
    \dfrac{1-(1+r)^{-t}}{r}, & r>0,\\[4pt]
    t, & r=0.
  \end{cases}
  \label{eq:annuity}
\end{equation}
Приведённые затраты на покупку и подписку соответственно составляют
\begin{align}
  C_{\mathrm{b}}(t)&=P-\frac{R(t)}{(1+r)^t},\\
  C_{\mathrm{s}}(t)&=sA(t,r).
\end{align}

При известном $t$ покупка предпочтительна, если $C_{\mathrm{b}}(t)<C_{\mathrm{s}}(t)$. При $R(t)=0$ и $r=0$ затраты равны при $t^*=P/s$. Для $R(t)=0$ и $r>0$ из равенства $P=sA(t,r)$ следует порог
\begin{equation}
  t^*=-\frac{\ln(1-rP/s)}{\ln(1+r)},\qquad rP<s.
  \label{eq:threshold}
\end{equation}
Первый целый срок, при котором покупка обходится не дороже подписки, равен $\lceil t^*\rceil$. При $R(t)=0$, $r>0$ и $rP\geq s$ подписка дешевле покупки при любом конечном сроке, поскольку $sA(t,r)<s/r\leq P$.

\subsection{Выбор при неопределённом сроке}

Потребитель, нейтральный к риску, выбирает вариант с меньшими ожидаемыми затратами. Из равенства $\mathbb E C_{\mathrm{b}}(T)=\mathbb E C_{\mathrm{s}}(T)$ следует порог абонентской платы
\begin{equation}
  s^*=
  \frac{P-\sum_i p_i R(t_i)(1+r)^{-t_i}}
       {\sum_i p_i A(t_i,r)}.
  \label{eq:price}
\end{equation}
Подписка предпочтительна при $s<s^*$, покупка~--- при $s>s^*$. При равенстве обе альтернативы эквивалентны по этому критерию.

Неприятие риска описывается полезностью $u(-C)=-\exp(\gamma C)$, где $\gamma>0$ измеряется в обратных денежных единицах. Максимизация $\sum_i p_i u(-C(t_i))$ эквивалентна минимизации денежного эквивалента затрат
\begin{equation}
  D_\gamma(C)=\frac{1}{\gamma}\ln\left(
      \sum_i p_i\exp\bigl(\gamma C(t_i)\bigr)\right).
  \label{eq:utility}
\end{equation}
При $\gamma\to0$ эта величина стремится к ожидаемым затратам. По неравенству Йенсена $D_\gamma(C)\geq\mathbb E C$, со строгим неравенством при непостоянных затратах. Если $R(t)=0$, стоимость покупки фиксирована, а стоимость подписки зависит от $T$. В этом случае учёт неприятия риска может изменить выбор в пользу покупки.

\subsection{Сценарный расчёт}

В условном примере $P=12\,000$~руб., $s=1\,000$~руб. в месяц, $r=0{,}01$ и $R(t)=0$. Вероятности и приведённые затраты по сценариям представлены в табл.~\ref{tab:scenarios}.

\begin{table}[htbp]
  \caption{Приведённые затраты по сценариям}
  \label{tab:scenarios}
  \centering
  \begin{tabular}{@{}rrrr@{}}
    \toprule
    $t_i$, мес. & $p_i$ & Покупка, руб. & Подписка, руб.\\
    \midrule
     6 & $0{,}30$ & $12\,000$ &  $5\,795$\\
    12 & $0{,}40$ & $12\,000$ & $11\,255$\\
    24 & $0{,}30$ & $12\,000$ & $21\,243$\\
    \bottomrule
  \end{tabular}
\end{table}

По формуле~\eqref{eq:threshold} $t^*\approx12{,}85$. При известном сроке покупка дешевле начиная с 13 месяцев. Для исходных вероятностей ожидаемые затраты подписки составляют $12\,614$~руб., а порог по формуле~\eqref{eq:price} равен $951{,}35$~руб. в месяц. При заданной цене подписки предпочтительна покупка.

Зафиксируем вероятность использования в течение 12 месяцев на уровне $0{,}40$. Обозначим вероятность 24 месяцев через $q$, тогда вероятность 6 месяцев равна $0{,}60-q$. Ожидаемые затраты подписки линейно возрастают с $q$ и сравниваются с ценой покупки при $q\approx0{,}2603$. Например, при $q=0{,}17$ ожидаемые затраты подписки составляют $10\,605$~руб. и подписка предпочтительна по критерию средних затрат. При том же $q$ и $\gamma=10^{-4}$~руб.$^{-1}$ денежный эквивалент затрат подписки по формуле~\eqref{eq:utility} составляет $12\,305$~руб. и по критерию ожидаемой полезности предпочтительна покупка.

Для эмпирической проверки потребуются сведения о тарифах, фактическом выборе и длительности использования продукта. Распределение сроков можно оценить по истории использования с учётом наблюдений, для которых пользование продуктом продолжалось к концу сбора данных. На отдельной выборке можно сравнить долю совпадений фактических решений с рекомендациями по двум критериям и оценить, улучшает ли учёт риска описание выбора.

\begin{thebibliography}{9}
\bibitem{bardhi}
F. Bardhi and G. M. Eckhardt, \textit{Access-Based Consumption: The Case of Car Sharing}, Journal of Consumer Research, vol.~39, no.~4, pp.~881--898, 2012. \url{https://doi.org/10.1086/666376}

\bibitem{lambrecht}
A. Lambrecht and B. Skiera, \textit{Paying Too Much and Being Happy About It: Existence, Causes, and Consequences of Tariff-Choice Biases}, Journal of Marketing Research, vol.~43, no.~2, pp.~212--223, 2006. \url{https://doi.org/10.1509/jmkr.43.2.212}

\bibitem{dellavigna}
S. DellaVigna and U. Malmendier, \textit{Paying Not to Go to the Gym}, American Economic Review, vol.~96, no.~3, pp.~694--719, 2006. \url{https://doi.org/10.1257/aer.96.3.694}

\bibitem{keeney}
R. L. Keeney and H. Raiffa, \textit{Decisions with Multiple Objectives: Preferences and Value Trade-Offs}. Cambridge University Press, 1993. \url{https://doi.org/10.1017/CBO9781139174084}
\end{thebibliography}
\end{document}
